\documentclass[UTF8]{ctexart}
\usepackage{amsmath,amssymb}
\usepackage[a4paper,margin=2.5cm]{geometry}
\title{数学分析 0921}
\date{}

\begin{document}
\maketitle

\subsection*{1. 实数的连续性}

\subsubsection*{Dedikind 连续性}

设 $S$ 为一个稠密有序数集。若把 $S$ 分成 $A,B$ 两类，满足

\begin{enumerate}
\item[(i)] \textbf{不空}：$A$ 与 $B$ 都至少包含 $S$ 中的一个数；
\item[(ii)] \textbf{不漏}：$S$ 中的每个数，或者属于 $A$，或者属于 $B$；
\item[(iii)] \textbf{不乱}：$A$ 中任一个数 $a$ 都小于 $B$ 中任一个数 $b$，即
$\forall a\in A,\ b\in B,\ a<b$。
\end{enumerate}

则称 $A,B$ 为 $S$ 的一个\textbf{分划}，记为 $A\mid B$；
$A$ 称为下类，$B$ 称为上类。

若 $S$ 的每个分划 $A\mid B$ 都存在唯一的 $c\in S$，使
$\forall a\in A,\ b\in B$ 有
\[
a\leq c\leq b,
\]
则称 $S$ 具有 Dedikind 连续性（即 $S$ 是连续的）。
这里的 $c$ 称为\textbf{分界点}。

上述定义的三个要点：

\begin{enumerate}
\item \textbf{稠密}：$\forall a,b\in S,\ a<b$，
$\exists c\in S,\ a<c<b$；
\item \textbf{存在且唯一}：$\exists!\,c\in S$，即分界点存在且唯一；
\item 分界点 $c\in S$：$A$ 中的数都不超过 $c$，
$B$ 中的数都不小于 $c$，即所有 $a\in A$、$b\in B$ 都被
$c$ 分在两边；而 $c$ 本身必须落在 $S$ 中，
要么在 $A$ 这一边，要么在 $B$ 这一边。
\end{enumerate}

由 $c$ 总属于 $A$ 或属于 $B$（“不漏”），
若 $c\in A$ 则 $c$ 是 $A$ 的最大数，若 $c\in B$ 则 $c$ 是 $B$ 的最小数，
所以连续性准则也可以这样说：若 $S$ 的每个分划 $A\mid B$，
或者 $A$ 有最大数，或者 $B$ 有最小数，则称 $S$ 是连续的。

\subsubsection*{$\mathbb Q$ 是不连续的}

找出 $\mathbb Q$ 的一个分划 $A\mid B$，
使该分划在 $\mathbb Q$ 内没有分界点即可。

考虑分划
\[
A=\{x\in\mathbb Q:x<\sqrt2\},\qquad
B=\{x\in\mathbb Q:x\geq\sqrt2\}.
\]
因为 $\sqrt2\notin\mathbb Q$，两集合构成 $\mathbb Q$
的一个分划（$A\cup B=\mathbb Q$ 且 $A\cap B=\varnothing$）。
下面证明 $A$ 中没有最大的数，同理 $B$ 中没有最小的数。

我们的目标是寻找一个 $\epsilon$，使得
\[
\forall a\in A,\ \exists\epsilon>0,\ \epsilon\in\mathbb Q,\quad
a+\epsilon\in A,
\]
即要求
\[
\begin{cases}
\epsilon>0,\\
\epsilon\in\mathbb Q,\\
a+\epsilon<\sqrt2.
\end{cases}
\]

先按放缩法估算。由 $a+\epsilon<\sqrt2$ 得
\[
(a+\epsilon)^2<2
\iff a^2+2a\epsilon+\epsilon^2<2.
\]
若 $\epsilon<2a$，则 $\epsilon^2<2a\epsilon$，于是
\[
a^2+4a\epsilon<2
\iff \epsilon<\frac{2-a^2}{4a}.
\]
故只需取
\[
\epsilon<\min\left\{2a,\ \frac{2-a^2}{4a}\right\},
\]
即可满足要求。
（注意：$\epsilon^2<2a\epsilon$ 与 $\epsilon<2a$ 只有在
$a>0$ 时才等价；而当 $a$ 靠近 $0$ 时，$2a$ 可能已经大于
$\sqrt2-a$，因此还须补上 $\epsilon<\sqrt2-a$ 这一条限制，
方能保证 $a+\epsilon<\sqrt2$，也才保证 $\epsilon>0$。）

下面给出一个更直接的取法。对任意 $a\in A$，
总有 $\sqrt2-a>0$，令
\[
\epsilon=\frac12\min\{1,\ \sqrt2-a\}>0 .
\]
由 $1,\sqrt2\in\mathbb Q$ 及 $a\in\mathbb Q$ 可知
$\epsilon\in\mathbb Q$；并且
\[
\epsilon\leq\frac12(\sqrt2-a)<\sqrt2-a,
\]
从而
\[
a+\epsilon<a+(\sqrt2-a)=\sqrt2,
\]
即 $a+\epsilon\in A$。
这说明 $A$ 中不存在最大的数：
对 $A$ 中任何一个数 $a$，总能在 $A$ 中找到一个比它更大的数。
（同理，对任意 $b\in B$，取
$\epsilon=\frac12\min\{1,\ b-\sqrt2\}$，
可证 $b-\epsilon\in B$，故 $B$ 中不存在最小的数。）

因此，$A$ 中没有最大的数，$B$ 中没有最小的数，
该分划在 $\mathbb Q$ 内没有分界点，
即不存在 $c\in\mathbb Q$ 使
$\forall a\in A,\ b\in B$ 有 $a\leq c\leq b$。
所以 $\mathbb Q$ 不具有 Dedikind 连续性，$\mathbb Q$ 不连续。

\subsubsection*{定理：实数集是 Dedikind 连续的}

连续需要证明三个条件：$S$ 是有序稠密集；存在一个分界点；
分界点要唯一。

\noindent\textbf{证：}

\noindent 1) 实数集显然有序且稠密。

\noindent 2) 存在分界点。我们的目标是依据分划的性质构造出一个数，
再证明它确实是分界点。

先回忆中学里讲过的，什么是实数。实数就是十进无穷小数
\[
a_0.a_1a_2a_3\cdots a_k\cdots,
\]
其中 $a_0$ 是整数，$a_k\ (k\geq1)$ 是满足
$0\leq a_k\leq9$ 的整数。

由 $A\cup B=\mathbb R$ 得 $\mathbb Z\subseteq A\cup B$，即
$\mathbb Z=(\mathbb Z\cap A)\cup(\mathbb Z\cap B)$；
又由 $A\cap B=\varnothing$ 得
$(\mathbb Z\cap A)\cap(\mathbb Z\cap B)=\varnothing$。
也就是说，分划点把整数分成两块：一块全在 $A$ 里，
一块全在 $B$ 里。

\textbf{声明 1：} $\exists a\in\mathbb Z$，使 $a\in A$ 且 $a+1\in B$。下证。

\textbf{声明 2：} $\exists i_0\in A$ 且 $i_0\in\mathbb Z$。

否则，若 $A$ 中一个整数都没有，即 $\mathbb Z\subseteq B$。
任取 $A$ 中的一个元素 $a$，它理应小于所有的整数，显然矛盾。

可以论证声明 1 了。若从 $i_0$ 起的整数都属于 $A$，
即 $\forall i\in\mathbb Z,\ i\geq i_0$ 都有 $i\in A$，
此时取定 $B$ 中任意数 $b$，则
$\forall i'\in\mathbb Z,\ i'\geq i_0$ 都有 $i'\leq b$；
但整数可以取得任意大，总能找到一个 $i'\geq i_0$ 且 $i'>b$，
这与“不乱”矛盾。所以在 $A$ 中一定可以找到一个最大整数，
任何大于它的整数都属于 $B$，即
\[
\exists a\in\mathbb Z,\quad a\in A\ \text{且}\ a+1\in B .
\]
那么我们想要构造的这个数，整数位就为 $i_0$。

对 $i_0\sim i_0+1$ 之间的区间作十等分，考虑
$i_0.0,\ i_0.1,\ \cdots,\ i_0.9$，其中必有落在 $A$ 中的数。

总存在一个 $i$，使得 $i_0.i$ 属于 $A$ 且
$i_0.(i+1)$ 属于 $B$。
否则，若对于 $\forall i_0.i\in A$ 都有
$i_0.(i+1)\in A$，一路推下去可得
$i_0.(i+10)\in A$，即 $(i_0+1).0\in A$；
而由上面整数位的构造，$i_0+1\in B$，与“不漏”矛盾。
当 $i=9$ 时发生进位，$i_0.(9+1)=(i_0+1).0$，
上面的写法已经把这一步包含在内。

我们就确定了这个数小数点后一位是 $i$，记它为 $i_1$。
如此类推，就可以构造出实数
\[
r=i_0.i_1i_2i_3\cdots .
\]

下面证明它确实是分界点。我们证明
$\forall a\in A,\ a\leq r$，另一边同理即可。

如果存在 $a_0\in A$ 使 $a_0>r$，把两个数都写成十进无穷小数，
则存在非负整数 $j_0$，使 $a_{j_0}>i_{j_0}$，
而当 $j<j_0$ 时 $a_j=i_j$。也就是说，$a_0$ 的第 $j_0$ 位上的数
不小于 $i_{j_0}+1$，因此
\[
a_0\geq i_0.i_1\cdots i_{j_0-1}(i_{j_0}+1)
>i_0.i_1\cdots i_{j_0-1}i_{j_0}=r .
\]
而由构造，$i_0.i_1\cdots i_{j_0-1}(i_{j_0}+1)\in B$。
于是 $a_0$ 不小于 $B$ 中的这个数，由“不乱”知 $a_0\notin A$，
与 $a_0\in A$ 矛盾。这就证明了
$\forall a\in A,\ a\leq r$。

同理可证 $\forall b\in B,\ r\leq b$。
所以 $r$ 确实是分划 $A\mid B$ 的分界点，存在性得证。

\noindent 3) 分界点唯一。

假设存在两个分界点且不相等，不妨设 $r_1<r_2$，取
\[
r=\frac{r_1+r_2}{2},
\]
显然 $r_1<r<r_2$。

对 $\forall a\in A$ 有 $a\leq r_1<r$，即
$\forall a\in A,\ a\neq r$。又因为 $r$ 也是分界点，
应有 $r\leq r_1$，与 $r_1<r$ 矛盾，所以 $r\notin A$。

同理对 $\forall b\in B$ 有 $r<r_2\leq b$，即
$\forall b\in B,\ b\neq r$；又因为 $r$ 也是分界点，
应有 $r_2\leq r$，与 $r<r_2$ 矛盾，所以 $r\notin B$。

但分划满足 $A\cup B=\mathbb R$，即 $r$ 总属于 $A$ 或属于 $B$，
矛盾。假设错误，分界点唯一。证毕。

\medskip

我们从自然数将数集扩充到了实数，使数集对于加减乘除都封闭，
且稠密并连续，提供了我们数学分析的舞台。

\subsection*{2. 集合、映射与函数}

映射：定义域，靶集，对应法则。

单射 injection/monomorphism；
满射 surjection/epimorphism；
双射 bijective/isomorphism。

\subsubsection*{映射的复合}

对于 $f:X\longrightarrow Y,\ g:W\longrightarrow Z$，
如果 $f(X)\subseteq W$，则定义 $g\circ f$ 是映射
\[
x\longmapsto g\bigl(f(x)\bigr).
\]

\textbf{四个命题：} $f,g$ 可以复合，
$f:X\longrightarrow Y$、$g:Y\longrightarrow Z$，
考虑 $g\circ f:X\longrightarrow Z$。

\begin{enumerate}
\item 如果 $f,g$ 都是单射，则 $g\circ f$ 也是单射；
\item 如果 $f,g$ 都是满射，则 $g\circ f$ 也是满射；
\item 如果 $g\circ f$ 是单射，则 $f$ 是单射；
\item 如果 $g\circ f$ 是满射，则 $g$ 是满射。
\end{enumerate}

\subsection*{3. 逻辑}

\end{document}
